% Section 2 -- Background (P2)
% Needs: \usepackage{amsmath, amssymb, booktabs}

\section{Background}
\label{sec:background}

This section introduces the three ideas the rest of the report builds on: Ising formulations of optimisation problems, the QAOA algorithm, and the cost of simulating quantum circuits on classical hardware.

\subsection{Ising models and QUBO}
\label{sec:ising}

Many combinatorial problems ask for the best assignment of $n$ yes/no decisions.
A common way to write them is as a \emph{quadratic unconstrained binary optimisation} (QUBO) problem: minimise $x^{\mathsf T} Q x$ over binary vectors $x \in \{0, 1\}^n$, for some matrix $Q$.
Substituting $x_i = (1 - z_i)/2$ turns a QUBO into an equivalent \emph{Ising model} over spins $z_i \in \{+1, -1\}$:
\begin{equation}
  C(z) = c_0 + \sum_i h_i z_i + \sum_{i<j} J_{ij}\, z_i z_j ,
  \label{eq:ising-general}
\end{equation}
where the $h_i$ are local fields, the $J_{ij}$ are couplings between pairs of spins, and $c_0$ is a constant.
Lucas~\cite{lucas2014} gives Ising formulations of many NP-hard problems, including number partitioning, Max-Cut and graph colouring.
Constraints are usually added as penalty terms that raise the cost of any assignment that breaks them, which is how we handle conflicting jobs in Section~\ref{sec:formulation}.

The Ising form maps directly onto qubits.
Replacing each spin $z_i$ by the Pauli operator $Z_i$ gives a Hamiltonian $H_C$ that is diagonal in the computational basis.
Its eigenvalue on the basis state $\lvert x \rangle$ is exactly the cost of the corresponding assignment, with $z_i = 1 - 2x_i$.
The ground state of $H_C$, its lowest-energy basis state, is therefore the optimal solution, and solving the problem means preparing and measuring a state close to it.

\subsection{The Quantum Approximate Optimization Algorithm}
\label{sec:qaoa-background}

QAOA~\cite{farhi2014} prepares a low-energy state of $H_C$ with a circuit of $p$ layers.
It starts from the uniform superposition $\lvert + \rangle^{\otimes n}$, in which every one of the $2^n$ assignments is equally likely.
Each layer applies two operations:
\begin{itemize}
  \item a \emph{cost layer} $U_C(\gamma) = e^{-i\gamma H_C}$, which multiplies each basis state by a phase that depends on its cost; and
  \item a \emph{mixer layer} $U_M(\beta) = e^{-i\beta H_M}$ with $H_M = \sum_i X_i$, which rotates every qubit about the $x$-axis and so mixes amplitude between assignments that differ by flipping one job.
\end{itemize}
The cost layer alone would not change any measurement probabilities, because it only changes phases.
The mixer turns those phase differences into amplitude differences, so that, with well-chosen angles, low-cost assignments become more likely.
With $p$ layers the circuit prepares
\begin{equation}
  \lvert \boldsymbol\gamma, \boldsymbol\beta \rangle
  = U_M(\beta_p)\, U_C(\gamma_p) \cdots U_M(\beta_1)\, U_C(\gamma_1)\, \lvert + \rangle^{\otimes n},
\end{equation}
which has $2p$ real parameters $\boldsymbol\gamma = (\gamma_1, \dots, \gamma_p)$ and $\boldsymbol\beta = (\beta_1, \dots, \beta_p)$.

QAOA is a \emph{hybrid} algorithm.
The quantum circuit estimates the energy $\langle H_C \rangle = \langle \boldsymbol\gamma, \boldsymbol\beta \rvert H_C \lvert \boldsymbol\gamma, \boldsymbol\beta \rangle$, and a classical optimiser adjusts the angles to lower it.
Once the angles are optimised, the circuit is measured many times, and each measurement gives one candidate schedule.

As $p$ grows, QAOA can imitate a slow adiabatic evolution, which is guaranteed to reach the ground state, so the result can only improve with depth when the angles are optimised well~\cite{farhi2014}.
At small $p$ the algorithm is approximate, and its quality depends on two practical issues.
First, deeper circuits contain more gates, which matters for both simulation cost and hardware noise.
Second, the energy as a function of the angles has many local minima, so the starting angles strongly affect the result.
Zhou et al.~\cite{zhou2020} proposed INTERP, which builds starting angles for depth $p + 1$ by interpolating the optimised angles from depth $p$.
Section~\ref{sec:method} describes how we use it.

\subsection{Simulating quantum circuits on GPUs}
\label{sec:simulation}

A classical statevector simulator stores the full quantum state of $n$ qubits as $2^n$ complex amplitudes.
Table~\ref{tab:memory} shows how fast this grows.
In single precision (two 32-bit floats per amplitude, 8 bytes in total), a 30-qubit state needs 8\,GiB, and every extra qubit doubles that.
Double precision doubles the memory again.

\begin{table}[h]
  \centering
  \caption{Memory needed to store an $n$-qubit statevector.}
  \label{tab:memory}
  \begin{tabular}{@{}lccccc@{}}
    \toprule
    Qubits $n$ & 20 & 24 & 28 & 30 & 32 \\
    \midrule
    Single precision (8 B per amplitude) & 8\,MiB & 128\,MiB & 2\,GiB & 8\,GiB & 32\,GiB \\
    Double precision (16 B per amplitude) & 16\,MiB & 256\,MiB & 4\,GiB & 16\,GiB & 64\,GiB \\
    \bottomrule
  \end{tabular}
\end{table}

Time grows just as fast.
Applying one gate updates every amplitude, so each gate costs work proportional to $2^n$.
A QAOA layer for our problem contains one $ZZ$ interaction for every pair of jobs, each built from three gates, plus one rotation per qubit: 590 gates per layer at $n = 20$.
Optimising the angles then needs hundreds of circuit evaluations.
This work consists of many simple, independent updates to a large array, which suits GPUs: they have thousands of cores and much higher memory bandwidth than CPUs.
CUDA-Q~\cite{cudaq} lets the same circuit run on a CPU simulator (target \texttt{qpp-cpu}) or a GPU simulator (target \texttt{nvidia}, built on NVIDIA's cuStateVec library) by changing one setting, which makes a direct comparison possible.

Two points about simulation are important for interpreting our results.
First, a GPU speeds up the \emph{simulation} of QAOA, not the algorithm itself.
The cost still doubles with every qubit, so a faster simulator adds only a few qubits to the largest problem we can study.
Second, simulating noise is far more expensive.
A noisy state is described by a density matrix with $4^n$ entries instead of $2^n$, so a 12-qubit noisy simulation stores as many numbers as a 24-qubit noiseless one.
This is why our noise study is limited to small problems.
